\subsection{Illustrative Planning Scenario and Reproduction by EXP-00}
\label{sec:practical-example}

Рассмотрим синтетический иллюстративный сценарий разработки небольшого сервиса
обработки заявок. Фиксированный
объём работы включает API-контракт, слой хранения, бизнес-логику, фоновую
обработку, наблюдаемость, тесты и упаковку. В условной шкале он равен
$38$ story points, или $19$ единицам $Z$; нормировка $X=4$ часа даёт
последовательный исходный срок $T_h=76$ часов. Эта связь используется только
для прозрачности сценария и не является эмпирическим преобразованием story
points в часы.

Task-DAG содержит последовательное ядро
$1\to2\to3\to4\to6$ и две боковые ветви: наблюдаемость с упаковкой и работы,
которые становятся доступными после фиксации интерфейсов. Сравниваются четыре
варианта организации работы. В вариантах 1--3 меняются режим и управляемый
параллелизм, поэтому меняются и коэффициенты $k_i,h_i$. Вариант 4 сохраняет
конфигурацию варианта 3, но дробит монолитную задачу тестирования на три части
с точными зависимостями. Только сравнение 3→4 выделяет эффект декомпозиции при
прочих равных.

\begin{table}[ht]
\centering
\caption{Сокращённое сравнение вариантов иллюстративного сценария}
\label{tab:compact-practical}
\begin{tabular}{lcccc}
\toprule
Показатель & Вариант 1 & Вариант 2 & Вариант 3 & Вариант 4 \\
\midrule
$P$ & 1 & 2 & 4 & 4 \\
Режим & интеракт. & смешанный & делегир. & делегир. \\
Декомпозиция & крупная & крупная & крупная & дроблёные тесты \\
$W_4$, ч & 65.0 & 59.8 & 53.6 & 53.6 \\
$W_4/P$, ч & 65.0 & 29.9 & 13.4 & 13.4 \\
$L_4$, ч & 51.8 & 47.8 & 43.2 & 36.7 \\
$H$, ч & 65.0 & 38.4 & 19.0 & 19.0 \\
$B_4$, ч & 65.0 & 47.8 & 43.2 & 36.7 \\
$T^*$, ч & 65.0 & 47.8 & 43.2 & 36.7 \\
Ускорение & $1.17\times$ & $1.59\times$ & $1.76\times$ & $2.07\times$ \\
\bottomrule
\end{tabular}
\end{table}

В серии EXP-00 для вариантов 2--4 точный решатель получил статус
\texttt{OPTIMAL}. В каждом
случае предъявлено допустимое расписание с $T(\pi)=B_4$, поэтому цепочка
$B_4\le T^*\le T(\pi)$ доказывает равенство $T^*=B_4$ именно для этих
сценариев. Равенство не предполагается универсальным свойством M4.

\subsubsection{Local Blocking and the Human Resource}

В варианте 2 два агентных слота обслуживаются одним разработчиком. Наивное
расписание, допустимое по DAG и числу агентов, содержит пересечения human-фаз
и потому недопустимо в M4. Исправленное расписание на
Рисунке~\ref{fig:practical-gantt} размещает постановки и проверки
последовательно. Задача 7 ожидает проверку $22.2$ часа и удерживает второй
слот, но ожидание идёт параллельно критической ветви и не увеличивает срок
проекта. Этот пример показывает, почему одна величина $Q(\pi)$ не определяет
makespan: важно положение ожидания в ресурсном расписании.
\begin{figure}[ht]
\centering
\includegraphics[width=\textwidth]{candidate_variant2_schedule_comparison.png}
\caption{Вариант 2 до и после явного учёта человеческого ресурса. Слева
исходное расписание содержит пересекающиеся human-фазы и недопустимо в M4.
Справа обслуживание разработчиком сериализовано; ожидающая задача удерживает
свой агентный слот, но не увеличивает общий срок.}
\label{fig:practical-gantt}
\end{figure}

\FloatBarrier

\subsubsection{Critical-Path Decomposition}

В варианте 3 вся задача тестирования ждёт завершения фонового обработчика. В
варианте 4 модульные и интеграционные тесты запускаются после собственных
предшественников, а в хвосте критического пути остаётся только проверка
повторных попыток. При неизменных $A=34.6$, $H=19.0$ и $W_4=53.6$ критический
путь и точный makespan сокращаются с $43.2$ до $36.7$ часа. Эффект равен
$6.5$ часа, или примерно $15\%$ срока варианта 3; ускорение относительно
исходного сценария возрастает с $1.76\times$ до $2.07\times$.

\begin{figure}[ht]
\centering
\includegraphics[width=\textwidth]{exp00_dag_variants_3_4.png}
\caption{Task-DAG и критический путь до и после уточнения зависимостей задачи
тестирования. Суммарная работа сохраняется, сокращается лишняя сериализация.}
\label{fig:practical-dag}
\end{figure}

\FloatBarrier

Полные параметры задач, фазовые интервалы, проверка человеческого ресурса и
элементарные примеры M0--M1 приведены в Приложении~\ref{app:examples}.
